\begin{abstract}
Software-engineering agents generate rich traces, including issue
interpretations, file searches, edits, test runs, failures, and recoveries,
that are discarded after each benchmark trial. Recycling these traces into
reusable text skills is an attractive way to let a frozen agent accumulate
experience, but it is also where most systems introduce leakage: raw trace
retrieval carries task-specific solutions into future prompts, and one-shot
``summarize the trajectory'' pipelines emit over-specific or actively harmful
instructions that no one audits. Either way, apparent gains come from
memorization or noise rather than transfer. We present \method{}, a
verifier-grounded framework that evolves external skills for \emph{frozen}
coding agents while making leakage structurally hard. \method{} refuses to call
every trace artifact a skill: it separates single-task evidence from reusable
skills, promotes evidence through a hierarchy of repo-level and cross-repo
behavior skills, and trains separate Skill Writer and Skill Evaluator
policies behind a fixed harness policy. The design is built around one hard
information boundary. During SWEGym training the official verifier supplies
outcome labels, evaluator calibration, and validation gates; at downstream
evaluation the verifier never touches the update path, so no hidden test, gold
patch, or future outcome may alter a skill. We instantiate the framework in a
runnable system (Pi harness, OpenAI-compatible models served from our own
self-hosted deployment, versioned skill archives with promotion logs and
validation gates) and specify an evaluation protocol over GLM-5.2 under two
agent harnesses, Pi and Claude-style, that isolates harness transfer, with
SWE-bench Verified and DeepSWE as in-domain targets, all under a shared
sandbox and acceptance rule.
\end{abstract}
